\section{Success Criteria}
For the adaptive controller, success is not defined solely by the minimum mean recovery time. A useful adaptive regime should keep the network functional under heterogeneity, preserve bounded recovery, and allow non-uniform recovery across agents when that reflects the state of the group.

\begin{itemize}
    \item The network remains operational under stress and returns to a recoverable regime.
    \item Recovery stays bounded, even if the mean recovery time is not minimal.
    \item Per-agent recovery is allowed to spread in time, provided the group does not collapse.
    \item The controller reacts to state, but does not turn into hard control.
\end{itemize}

In other words, a good adaptive outcome is not necessarily the fastest one. It is the one that supports live, heterogeneous, and stable group operation.

\selectlanguage{russian}
\section{Критерии успеха}
Для адаптивного контроллера успех не сводится только к минимальному среднему времени восстановления. Полезный режим должен сохранять работоспособность сети при неоднородности, удерживать восстановление в конечных пределах и допускать неравномерное восстановление отдельных агентов, если это соответствует состоянию группы.

\begin{itemize}
    \item Сеть остается работоспособной под стрессом и возвращается в восстановимый режим.
    \item Восстановление остается конечным, даже если среднее время не минимально.
    \item Восстановление по агентам может растягиваться во времени, если группа не рушится.
    \item Контроллер реагирует на состояние, но не превращается в жесткое управление.
\end{itemize}

Иными словами, хороший адаптивный режим не обязательно самый быстрый. Он тот, который поддерживает живую, неоднородную и устойчивую работу группы.
\selectlanguage{english}
